\documentclass{situationist-preprint}

\title{Benatar's Asymmetry Again}
\msid{EXI-12-b}
\msfamily{EXI-12}
\msdiscipline{Meaning, Death and the Absurd}
\mstype{Consequence}
\mskeywords{Benatar; asymmetry; procreation; anti-natalism; complaints}

\begin{document}
\maketitle

\begin{abstract}
ETH-20-a shows that Benatar's asymmetry between the absence of pain and the absence of pleasure fails, and EXI-12-a concludes that a life containing more good than bad is better than non-existence. Benatar's case for the asymmetry, however, is partly explanatory: he argues that it alone explains four other asymmetries in our judgments about procreation, regret and absent people. We derive from the two results that these four asymmetries have a different explanation, an asymmetry of complaints: a wrong to a person requires someone who has a complaint, and the absence of pleasure in a world without her is bad for no one though it is not neutral. The explanatory argument therefore gives no support to the asymmetry.
\end{abstract}

\section{The explanatory argument}

\citet[ch.~2]{Benatar2006} supports his basic asymmetry, that the absence of pain is good even if no one enjoys it while the absence of pleasure is not bad unless someone is deprived, by the claim that it best explains four further asymmetries. (i) There is a duty not to create a person whose life would be miserable, but no duty to create a person whose life would be happy. (ii) It is odd to cite the benefit to a child as a reason for creating it, but not odd to cite its suffering as a reason against. (iii) We may regret, for the child's sake, having created a suffering child, but not, for anyone's sake, having failed to create a happy one. (iv) We are sad about suffering people in distant places, but not about the absence of happy people on an uninhabited island.

If the basic asymmetry fails, as Corollary~3.4 of ETH-20-a shows, then either the four asymmetries are unexplained or they have another explanation. We give the other explanation and show that it is consistent with the results of ETH-20 and EXI-12.

\section{The asymmetry of complaints}

\begin{principle}[Complaints]\label{p:compl}
An act wrongs a person only if it is worse for her or sets back an interest of hers; reasons to bring about an impersonally better outcome that is better for no one are not owed to anyone, and do not by themselves ground a duty.
\end{principle}

The principle is a familiar person-affecting view of wronging \citep{McMahan2009}. It is compatible with ETH-20-a, which holds that the absence of pleasure is bad for no one but impersonally worse than its presence, and with Principle~P1 of EXI-12-a, which compares outcomes impersonally.

\begin{proposition}\label{p:explain}
Principle~\ref{p:compl}, together with Corollary~3.4 of ETH-20-a and Theorem~3.1 of EXI-12-a, explains asymmetries (i)--(iv) without Benatar's basic asymmetry.
\end{proposition}

\begin{proof}
(i) The miserable person, once created, has a life that is bad for her and so has a complaint; by Principle~\ref{p:compl} creating her can wrong her. The happy person who is never created does not exist, so no one has a complaint, and the impersonal reason to create her, which exists by Corollary~3.4 of ETH-20-a, grounds no duty. (ii) Citing the child's benefit is not odd as a reason: by Theorem~3.1 of EXI-12-a a good life is a non-comparative benefit to the one who lives it. What is odd is to treat it as a reason \emph{owed} to the child, since there was no one to whom it was owed. (iii) Regret for the child's sake requires a child for whose sake one regrets; the never-created happy person is no one, so the regret can only be impersonal, and is weak, as the reason was. (iv) Sadness is fitting at what is bad for someone. The absent pleasures of the island are bad for no one (ETH-20-a), so personal sadness is unfitting, although the island world is impersonally worse than one containing happy lives.
\end{proof}

The explanation draws on what ETH-20-a concedes to Benatar, that the absence of pleasure is bad for no one, and not on what it denies, that this absence is neutral. The four asymmetries are asymmetries of complaint and of person-directed attitude; Benatar's basic asymmetry is an asymmetry of value. Only the first is needed to explain them.

\begin{objection}[Procreative duty]
If the absence of pleasure is impersonally bad, there is a reason to create every happy person, and the view implies a duty to procreate.
\end{objection}

\begin{reply}
By Principle~\ref{p:compl}, an impersonal reason that is owed to no one does not by itself ground a duty. Creating a happy person is permissible, as ETH-20-a and ETH-20-b conclude, and good, as EXI-12-a concludes, without being required.
\end{reply}

\section{Conclusion}

Benatar's four asymmetries are real, but they are explained by an asymmetry of complaints that leaves the comparison of EXI-12-a intact. The explanatory argument gives no independent support to the basic asymmetry, and a life with more good than bad remains better than never having been.

\begin{thebibliography}{9}\small
\bibitem[Benatar(2006)]{Benatar2006} Benatar, D. (2006). \emph{Better Never to Have Been: The Harm of Coming into Existence}. Oxford University Press.
\bibitem[McMahan(2009)]{McMahan2009} McMahan, J. (2009). Asymmetries in the morality of causing people to exist. In M.~A. Roberts and D.~T. Wasserman (eds.), \emph{Harming Future Persons}. Springer.
\end{thebibliography}

\end{document}
